\section{Tổng kết}

Dựa trên danh mục 13 chức năng và kế hoạch phân công ở Mục 3.2, nhóm tiến hành đối chiếu với \textit{``Quy định đồ án cuối kỳ''} môn học \textit{Vật lý cho Công nghệ Thông tin}:
\begin{equation*}
\text{Điểm CK} = (\text{Điểm YCCB} + \text{Điểm YCNC} + \text{Điểm BC}) \times K
\end{equation*}
Trong đó: Điểm YCCB tối đa 5.0đ (4.0đ luồng căn bản + 1.0đ tích hợp); Điểm YCNC tối đa 4.0đ (đăng ký các yêu cầu với phần cứng đôi một khác nhau); Điểm BC tối đa 1.0đ; Hệ số $K \in [0\%, 100\%]$ qua vấn đáp; Điểm cộng sản phẩm thật tối đa +1.0đ khi đóng gói hộp in 3D vận hành ổn định.

\subsection{Bảng tổng kết điểm số dự kiến của các thành viên}

\begin{center}
\small
\renewcommand{\arraystretch}{1.02}
\begin{tabularx}{\textwidth}{|c|X|c|c|c|c|}
\hline
\textbf{STT} & \textbf{Họ và tên thành viên} & \textbf{\parbox{2.2cm}{\centering Điểm căn bản\\(Tối đa 5.0đ)*}} & \textbf{\parbox{2.2cm}{\centering Điểm nâng cao\\(Tối đa 4.0đ)**}} & \textbf{\parbox{2.1cm}{\centering Điểm báo cáo\\(Tối đa 1.0đ)***}} & \textbf{\parbox{1.8cm}{\centering Tổng điểm\\dự kiến}} \\ \hline
1 & \textbf{Trần Nguyễn Tuấn Minh} \newline \footnotesize MSSV: 24127084 & 5.0 / 5.0 & 4.0 / 4.0 & 1.0 / 1.0 & \textbf{10.0} / 10.0 \\ \hline
2 & \textbf{Nguyễn Quốc Dương} \newline \footnotesize MSSV: 24127349 & 5.0 / 5.0 & 4.0 / 4.0 & 1.0 / 1.0 & \textbf{10.0} / 10.0 \\ \hline
3 & \textbf{La Anh Hào} \newline \footnotesize MSSV: 24127361 & 5.0 / 5.0 & 4.0 / 4.0 & 1.0 / 1.0 & \textbf{10.0} / 10.0 \\ \hline
\end{tabularx}
\captionof{table}{Bảng tổng kết điểm số dự kiến của từng thành viên theo Barem quy định}
\label{tab:tong_ket_diem}
\end{center}

\subsection{Chi tiết phân bổ tiêu chí Barem môn học cho từng thành viên}

\begin{center}
\scriptsize
\renewcommand{\arraystretch}{0.92}
\begin{tabularx}{\textwidth}{|p{2.5cm}|c|X|c|c|}
\hline
\textbf{Thành viên} & \textbf{Mã Barem} & \textbf{Nội dung triển khai thực tế trong hệ thống} & \textbf{Khung điểm} & \textbf{Đăng ký} \\ \hline
\multirow{4}{=}{\textbf{Trần Nguyễn Tuấn Minh} \newline \tiny MSSV: 24127084} 
 & YCCB & Luồng Input $\rightarrow$ ESP32 $\rightarrow$ MQTT $\rightarrow$ Web (DHT22) + Tích hợp & 4.0 + 1.0 & 5.0 / 5.0 \\ \cline{2-5}
 & YC4 & Lưu dữ liệu cảm biến/lịch sử theo thời gian lên Supabase PostgreSQL & 1.0 -- 1.5 & 1.5 / 1.5 \\ \cline{2-5}
 & YC5 & Lấy dữ liệu Cloud hiển thị Web dạng Gauge và Chart (Node-RED) & 0.5 -- 1.5 & 1.5 / 1.5 \\ \cline{2-5}
 & YC11 & Tự xây dựng mô hình AI phân tích điều kiện học tập tối ưu & 1.0 -- 2.0 & 2.0 / 2.0 \\ \hline

\multirow{5}{=}{\textbf{Nguyễn Quốc Dương} \newline \tiny MSSV: 24127349}
 & YCCB & Luồng Input $\rightarrow$ ESP32 $\rightarrow$ MQTT $\rightarrow$ Web (RFID RC522) + Tích hợp & 4.0 + 1.0 & 5.0 / 5.0 \\ \cline{2-5}
 & YC1 & Luồng Input $\rightarrow$ ESP32 $\rightarrow$ Output (Nút nhấn $\rightarrow$ LED đỏ 5mm) & 0.5 -- 1.5 & 1.5 / 1.5 \\ \cline{2-5}
 & YC2 & Luồng Input $\rightarrow$ Web (Quang trở LM393 $\rightarrow$ ESP32 $\rightarrow$ MQTT $\rightarrow$ Web) & 1.0 -- 1.5 & 1.5 / 1.5 \\ \cline{2-5}
 & YC6 & Thông báo về điện thoại khi học quá giờ bằng Telegram Bot API & 1.0 & 1.0 / 1.0 \\ \cline{2-5}
 & YC7 & Thông báo cho người dùng bằng Email tổng kết tuần qua Gmail SMTP & 1.0 & 1.0 / 1.0 \\ \hline

\multirow{4}{=}{\textbf{La Anh Hào} \newline \tiny MSSV: 24127361}
 & YCCB & Luồng Web $\rightarrow$ MQTT $\rightarrow$ ESP32 $\rightarrow$ Output (Active Buzzer 5V) + Tích hợp & 4.0 + 1.0 & 5.0 / 5.0 \\ \cline{2-5}
 & YC3 & Luồng Web $\rightarrow$ ESP32 $\rightarrow$ Output (Web $\rightarrow$ ESP32 $\rightarrow$ Màn hình LCD 1602) & 1.0 -- 1.5 & 1.5 / 1.5 \\ \cline{2-5}
 & YC9 & Bảo mật hệ thống: Đăng ký \& đăng nhập lưu tài khoản trong database & 1.0 -- 1.5 & 1.5 / 1.5 \\ \cline{2-5}
 & YC12 & [Sản phẩm thật] Cấu hình Wi-Fi mới qua Web Server SoftAP trên ESP32 & 1.0 -- 2.0 & 2.0 / 2.0 \\ \hline

\multicolumn{4}{|l|}{\textbf{Điểm Báo cáo cuối kỳ (Nộp đầy đủ, đúng hạn quy định trên Moodle):}} & \textbf{1.0 / 1.0} \\ \hline
\end{tabularx}
\captionof{table}{Bảng chi tiết đối chiếu các tiêu chí Barem môn học cho từng thành viên}
\label{tab:chi_tiet_barem}
\end{center}

\clearpage
\subsection{Minh chứng tính hợp lệ và tuân thủ Barem môn học}
\begin{enumerate}[leftmargin=1.0cm, noitemsep, topsep=1pt]
    \item \textbf{Tính đôi một khác nhau của thiết bị phần cứng:} 7 thiết bị ngoại vi sử dụng trong hệ thống gồm: DHT22, RFID RC522, Active Buzzer, LM393, Nút nhấn cơ học, LED đỏ 5mm, LCD 1602 I2C là \textbf{hoàn toàn độc lập và đôi một khác nhau}, thỏa mãn triệt để quy định tại Mục C.1 và C.2 của Barem.
    \item \textbf{Tính độc lập của từng yêu cầu nâng cao:} Mỗi yêu cầu trong Barem chỉ do duy nhất 1 thành viên phụ trách (Minh: YC4, YC5, YC11; Dương: YC1, YC2, YC6, YC7; Hào: YC3, YC9, YC12). Không có bất kỳ sự trùng lặp nào.
    \item \textbf{Tính tách biệt giữa dịch vụ thông báo nhanh và email:} Sử dụng Telegram Bot API cho thông báo nhanh tức thì (YC6) và sử dụng độc lập giao thức Gmail SMTP cho thư báo cáo tuần (YC7).
\end{enumerate}

\subsection{Cam kết thực hiện đồ án và liêm chính học thuật}

Nhóm thực hiện đề tài \textbf{Trạm học tập Pomodoro (Nhóm 07, Lớp 25C07)} xin cam kết:
\begin{enumerate}[leftmargin=1.0cm, noitemsep, topsep=1pt]
    \item \textbf{Sản phẩm thật 100\%:} Chế tạo và hoàn thiện phần cứng thực tế đóng gói hoàn chỉnh trong vỏ hộp in 3D, kết nối ổn định với máy chủ Node-RED và Supabase Cloud theo đúng các mốc tiến độ đã lập.
    \item \textbf{Tính ổn định chức năng:} Hiện thực hóa đầy đủ 13 chức năng đã thuyết minh; cam kết không thay đổi quá 1 chức năng theo đúng quy chế môn học.
    \item \textbf{Liêm chính học thuật:} Toàn bộ sơ đồ, mã nguồn, bản vẽ kỹ thuật và văn bản thuyết minh đều do các thành viên tự nghiên cứu, thiết kế và phát triển; tuân thủ nghiêm ngặt quy định về liêm chính học thuật.
\end{enumerate}

\begin{thebibliography}{99}
    \setlength{\itemsep}{0pt}
    \setlength{\parskip}{0pt}
    \small
    \bibitem{espressif2023} Espressif Systems, \textit{ESP32 Series Datasheet}, Version 4.1, 2023.
    \bibitem{mqttspec} OASIS Standard, \textit{MQTT Version 3.1.1 Specification}, 2014.
    \bibitem{nodered2024} OpenJS Foundation, \textit{Node-RED: Low-code programming for event-driven applications}, Official Documentation, 2026.
    \bibitem{supabase2024} Supabase Inc., \textit{The Open Source Firebase Alternative Documentation}, 2026.
    \bibitem{mfrc522} NXP Semiconductors, \textit{MFRC522 Standard Contactless Reader IC Datasheet}, Rev. 3.9, 2016.
    \bibitem{dht22} Aosong Electronics, \textit{Digital-output relative humidity \& temperature sensor/module DHT22}, 2015.
    \bibitem{fusion360} Autodesk Inc., \textit{Autodesk Fusion 360 User Guide and CAD Specifications}, 2026.
\end{thebibliography}
